% BHU PYQ -- Manually transcribed & error-corrected
\documentclass[11pt,a4paper]{article}

\usepackage[a4paper,top=2.5cm,bottom=2.5cm,left=2.8cm,right=2.8cm,headheight=14pt]{geometry}
\usepackage{amsmath,amssymb}
\usepackage[T1]{fontenc}
\usepackage[utf8]{inputenc}
\usepackage{lmodern}
\usepackage{microtype}
\usepackage{fancyhdr}
\usepackage{enumitem}
\usepackage{hyperref}
\usepackage{lastpage}
\usepackage{parskip}

\hypersetup{colorlinks=true,urlcolor=black,linkcolor=black,
  pdftitle={CHB-201: Inorganic & Physical Chemistry-I -- BHU PYQ},pdfauthor={Banaras Hindu University}}

\pagestyle{fancy}\fancyhf{}
\renewcommand{\headrulewidth}{0.4pt}\renewcommand{\footrulewidth}{0.4pt}
\fancyhead[L]{\small\textit{Banaras Hindu University}}
\fancyhead[R]{\small\textit{CHB-201: Inorganic \& Physical Chemistry-I}}
\fancyfoot[C]{\small Page \thepage\ of \pageref{LastPage}}

\newlist{parts}{enumerate}{2}
\setlist[parts,1]{label=(\alph*),leftmargin=*,itemsep=5pt,topsep=3pt}
\setlist[parts,2]{label=(\roman*),leftmargin=*,itemsep=3pt,topsep=2pt}
\newcommand{\pts}[1]{\hfill\textit{[#1]}}

\begin{document}

\begin{center}
  {\large\bfseries BANARAS HINDU UNIVERSITY}\\[2pt]
  {\normalsize B.Sc. (Hons.) Semester II Examination, 2023-24}\\[6pt]
  \rule{0.8\linewidth}{0.5pt}\\[6pt]
  {\large\bfseries Paper No. CHB-201: Inorganic \& Physical Chemistry-I}\\[4pt]
  {\normalsize Time: 3 Hours\hfill Maximum Marks: 70}
\end{center}
\rule{\linewidth}{0.4pt}
\small
\noindent\textit{Instructions: Answer five questions in total, including Question~1 from each section which is compulsory. Use separate answer books for Section~A and Section~B.}
\normalsize
\bigskip

\begin{center}
  {\large\bfseries SECTION A}\\[2pt]
  {\normalsize (Inorganic Chemistry-I -- Marks: 35)}
\end{center}
\rule{0.4\linewidth}{0.2pt}
\smallskip

\noindent\textbf{Question 1.} Answer all parts of the following: \pts{5 $\times$ 2 = 10}
\begin{parts}
  \item Differentiate between a nitride and an azide.
  \item What do you mean by sodide/potasside or alkalide ions in general?
  \item Compare $\text{BH}_3$ and $\text{CH}_4$ in terms of total number of available orbitals and electron pairs in their valence shells.
  \item Comment on the $2\text{e}^-\text{-}3\text{c}$ bonds of diborane and write one preparation method.
  \item Write the names of the hydrides of the nitrogen family and discuss their relative stability and reducing power.
\end{parts}

\medskip\rule{\linewidth}{0.2pt}

\noindent\textbf{Question 2.}
\begin{parts}
  \item Predict hybridization, draw structures, and explain shapes of: (i) $\text{XeF}_4$ and (ii) $\text{XeOF}_4$. \pts{6}
  \item Describe the structure of basic beryllium acetate. Explain why beryllium forms this basic acetate but other alkaline earth metals do not. \pts{6}
\end{parts}

\medskip\rule{\linewidth}{0.2pt}

\noindent\textbf{Question 3.}
\begin{parts}
  \item Discuss the interstitial and covalent carbides, giving their preparations and reactions. \pts{6}
  \item Discuss graphite intercalation compounds (GICs). Why is graphite conductive whereas diamond is an insulator? \dots \pts{6}
\end{parts}

\medskip\rule{\linewidth}{0.2pt}

\noindent\textbf{Question 4.}
\begin{parts}
  \item State the Kapustinskii equation for lattice energy and discuss how it is useful for calculating the lattice energy of crystals where ionic radii are not exactly known. \pts{6}
  \item Discuss the structure of crown ethers and cryptands. Explain how they are used for alkali metal complexation. \pts{6}
\end{parts}

\newpage
\begin{center}
  {\large\bfseries SECTION B}\\[2pt]
  {\normalsize (Physical Chemistry-I -- Marks: 35)}
\end{center}
\rule{0.4\linewidth}{0.2pt}
\smallskip

\noindent\textbf{Question 5.} Answer all parts of the following: \pts{5 $\times$ 2 = 10}
\begin{parts}
  \item What is the unit of $R$ (universal gas constant) in SI and CGS units?
  \item Arrange in decreasing order of magnitude: root-mean-square velocity, average velocity, and most probable velocity.
  \item Explain the Joule-Thomson experiment and find the value of the Joule-Thomson coefficient ($\mu_{JT}$) for an ideal gas.
  \item Define and derive the expression for critical temperature ($T_c$) of a van der Waals gas.
  \item Justify the statement: ``To maintain a small bubble needs excess pressure than a bigger one.''
\end{parts}

\medskip\rule{\linewidth}{0.2pt}

\noindent\textbf{Question 6.}
\begin{parts}
  \item Explain the statement: ``A gas can be liquefied at $T < T_c$ and $P > P_c$.'' Discuss the liquefaction of gases. \pts{6}
  \item Sketch the Carnot cycle in an $S$-$V$ (Entropy-Volume) diagram, labelling all the states involved. \pts{6}
\end{parts}

\medskip\rule{\linewidth}{0.2pt}

\noindent\textbf{Question 7.}
\begin{parts}
  \item What is a pseudo-first-order reaction? Give one example and write the expression for its rate constant. \pts{6}
  \item Derive the rate equation for a second-order reaction of the type $2\text{A} \to \text{Products}$. Calculate its half-life period. \pts{6}
\end{parts}

\medskip\rule{\linewidth}{0.2pt}

\noindent\textbf{Question 8.}
\begin{parts}
  \item Define surface tension. Discuss the drop-weight method for the determination of the surface tension of a liquid. \pts{6}
  \item Write short notes on: \pts{6}
  \begin{parts}
    \item Viscosity and its temperature dependence
    \item Maxwell-Boltzmann distribution of molecular speeds
  \end{parts}
\end{parts}

\vfill
\rule{\linewidth}{0.4pt}
\begin{center}\small
  \href{https://bhu-exam.vercel.app/}{Click here to visit our website}\\[4pt]
  \href{https://me-aryan.vercel.app/}{Click here to visit our contributor}\\[4pt]
  \href{https://quashan-me.vercel.app/}{Click here to visit our contributor}
\end{center}

\end{document}
